\subsection{Elastic Deformation of Rocks}

In the mechanics of solids, deformation means a change in the position, shape, or volume of a body under the action of external loading or internal force. In geological materials, deformation may be caused by burial, tectonic loading, pore-pressure changes, thermal expansion or contraction, and transient dynamic disturbances such as seismic or thermoelastic waves. From the geomechanical point of view, rocks must be considered together with their pore space and discontinuities, because natural rock masses commonly include fractures, faults, joints, and other structural features that influence how forces are transmitted and how deformation is distributed. For this reason, deformation in rocks is not only a geometric change of form, but also a physical response controlled by lithology, fabric, porosity, crack density, and stress history \cite{jaeger2007,cornet2007}.

Elastic deformation is the reversible part of deformation. When a body is loaded within its elastic range, it changes shape or volume, but it can return approximately to its original configuration after the load is removed. This reversibility is important because elastic deformation stores mechanical energy, and that stored energy can later be released. In rocks, elastic behavior is therefore the mechanical basis of wave propagation: a local disturbance produces stress and strain, and the tendency of the material to recover its original state provides the restoring mechanism that allows the disturbance to travel through the medium. Within the linear elastic approximation, stress is taken to be proportional to strain, which is the classical content of Hooke's law in continuum mechanics and seismology \cite{aki2002,jaeger2007}.

Elastic deformation must be distinguished from permanent or failure-related modes of rock response. If deformation remains after unloading, the behavior is no longer purely elastic. In mechanics this irreversible response is called plastic deformation, while in geological language permanent distributed deformation without macroscopic rupture is commonly described as ductile. Brittle deformation, by contrast, is characterized by cracking, fracture, or other forms of localized failure. Rocks can pass through these regimes progressively: they may respond elastically at small load, then accumulate irreversible strain, and finally fracture if the imposed stress continues to increase. Which regime dominates depends on lithology, temperature, confining pressure, fluid conditions, and the rate at which stress is applied \cite{jaeger2007,cornet2007}.

For wave propagation, the assumption of small elastic deformation is especially important. A propagating wave does not usually require the rock to move far from its initial equilibrium state. Instead, the wave produces a small perturbation of displacement, strain, and stress superimposed on the pre-existing geological stress field. Under such small-amplitude conditions, it is usually reasonable to approximate the rock as linearly elastic, even if the same rock would exhibit nonlinear, plastic, or brittle behavior under larger loading. In this sense, wave theory often deals with incremental deformation around a reference state rather than with the full long-term rheological behavior of the rock \cite{mavko2009,aki2002}.

This approximation is also justified because the equations of elastic-wave motion are formulated most naturally for infinitesimal strains. In the classical theory, the displacement field is assumed to vary smoothly enough that only first-order displacement gradients need to be retained. Once the strains are small, rigid-body translation and rotation can be separated from true distortion, and the stress--strain relation may be treated in linear form. This makes the mechanical description of wave propagation clear and physically interpretable, while still capturing the main dependence of wave speed and polarization on density and elastic stiffness \cite{aki2002,mavko2009}.

Nevertheless, the linear elastic assumption has definite limitations in real geological media. Many rocks are porous, microcracked, fractured, weathered, or compositionally heterogeneous. In such materials, elastic moduli may depend on confining stress, pore pressure, crack closure, and loading amplitude. Compliant crack-like parts of the pore space, including microcracks and compliant grain boundaries, can have a disproportionately large effect on elastic response. This means that the apparent stiffness measured during very small dynamic disturbances may differ from that observed in larger quasi-static loading, especially when cracks open, close, or slide \cite{mavko2009,jaeger2007}.

A further limitation arises from the geological structure of rock masses. Continuum elasticity treats the medium as if material properties vary smoothly from point to point, but real rock masses often contain bedding surfaces, foliation, joints, veins, and faults. The continuum approximation becomes meaningful only when an equivalent representative volume can be defined. At sufficiently large scale this simplification is very useful, but near discrete fractures, damage zones, or strongly layered intervals the actual deformation may become highly nonuniform. Under such conditions, stress concentration, crack-tip effects, local slip, and anisotropic compliance may significantly modify the response to thermal or mechanical disturbance \cite{cornet2007,jaeger2007,schultz2019}.

For the purposes of thermoelastic wave propagation, elastic deformation should therefore be understood as a first-order approximation that is both physically meaningful and geologically limited. It is meaningful because rocks can support small reversible stress and strain disturbances, and those disturbances constitute elastic waves. It is limited because real geological media are rarely ideal solids: their response may deviate from perfect linear elasticity when pore structure, fractures, mineral mismatch, damage, or strong thermal stressing become important. The following section develops the basic language of stress and strain that makes this elastic description possible \cite{jaeger2007,mavko2009,cornet2007}.